\documentclass[12pt,a4paper]{article}

\makeatletter
\def\input@path{{/home/maenwe/Master_CSM/modèle_cours_latex}}
\makeatother

\usepackage{prise_de_note}
\graphicspath{{images}}
\begin{document}
	\pagenumbering{gobble}
	\begin{titlepage}
		\centering
		
		% Logo
		\begin{figure}[h]
			\centering
			% Ajustez la largeur si besoin (ex: 0.4 ou 0.6)
			\includegraphics[width=0.5\textwidth]{/home/maenwe/Master_CSM/modèle_cours_latex/logo.jpg} 
		\end{figure}
		
		\vspace{2cm}
		
		% Filet horizontal du haut
		\rule{\linewidth}{0.5mm} \\[0.4cm]
		
		% Titre
		{\color{bleuTitre} \Huge \bfseries Estimation de paramètres et optimisations} \\[0.4cm]
		
		% Filet horizontal du bas
		\rule{\linewidth}{0.5mm}
		
		\vspace{1.5cm}
		
		{\Large \textbf{TD2}}
		
		\vfill
		
		% Bloc Informations
		\begin{minipage}{0.45\textwidth}
			\begin{flushleft} \large
				\textbf{Enseignants :} \\
				M. Abdelaziz \textsc{BELMILOUDI}
			\end{flushleft}
		\end{minipage}
		\hfill
		\begin{minipage}{0.45\textwidth}
			\begin{flushright} \large
				\textbf{Étudiant :} \\
				Nathan \textsc{HASCOET}
			\end{flushright}
		\end{minipage}
		
		\vspace{2cm}
		
		% Date
		{\large \today}
		
	\end{titlepage}
	
	\newpage
	\pagenumbering{arabic}
	\setcounter{page}{1}
	
	
	\begin{sol}
		$ $\\
		\begin{minipage}[t]{17cm}
			\begin{enumerate}
				\item l'application $\varphi : v \mapsto \norme{Av-b}$ est continue, de plus, en choisissant $v_0 \in \R^n$, il suit $\inf_v \norme{Av-b} \leq \norme{Av_0-b}$. Ainsi $\inf_v \norme{Av-b} \in \varphi^{-1}([0,\,\norme{Av_0-b}])$, de plus $\varphi^{-1}([0,\,\norme{Av_0-b}])$ est compact. Donc il existe $v \in \R^n$ tel que $\varphi(v) = \inf_v \norme{Av-b}$. De plus un tel $v$ est unique si le rang de $A$ est maximal : $\mathrm{rang}(A) = m$.
				
				\item Puisque $\varphi^2$ est convexe et différentiable : $v \in \mathrm{Argmin} \varphi(v)^2 \iff 0 = \nabla \varphi^2(v) = 2A^T(Av-b)$.\\
				Ainsi le problème admet une unique solution si et seulement si $A^TA$ est inversible, ce qui revient à demander à ce que $A^TA$ est définie positive car elle est symétrique réelle. Or pour tout $v \in \R^n$, $\left<A^TAv,v\right> = 0 \iff \left<Av,Av\right> = 0 \iff Av = 0$. Ainsi $A^TA$ est définie positive ssi $A$ est injective, ie. $\mathrm{rang}(A) = m$.
				
				\item 
			\end{enumerate}
		\end{minipage}
	\end{sol}
	
\end{document}